\chapter{FUNDAMENTAÇÃO TEÓRICA}
\label{cap:fundamentacao_teorica}

Esta seção revisa os principais conceitos teóricos que sustentam este trabalho, englobando as arquiteturas de aprendizado profundo, a formulação do fluxo óptico, técnicas de pré-processamento, métricas de avaliação e trabalhos relacionados.

\section{Segmentação Semântica e a Arquitetura U-Net}
\label{sec:unet}

A segmentação semântica consiste na tarefa de classificar cada pixel de uma imagem digital em uma categoria semântica predefinida. No contexto do Aprendizado Profundo (*Deep Learning*), as Redes Neurais Convolucionais Baseadas em Codificador-Decodificador (*Encoder-Decoder*) tornaram-se o estado da arte para essa finalidade. Dentre essas topologias, a U-Net destaca-se por sua capacidade de operar eficientemente com conjuntos de dados relativamente restritos \cite{ronneberger2015unet}.

A arquitetura da U-Net é simétrica e divide-se em duas macroestruturas:
\begin{enumerate}
    \item \textbf{Caminho de Contração (Codificador):} Funciona como uma rede convolucional típica, aplicando sucessivas camadas de convolução seguidas por ativações não lineares (como ReLU) e operações de subamostragem (*Max-Pooling*). Esse processo reduz progressivamente a resolução espacial da imagem enquanto extrai mapas de características (*feature maps*) de alto nível semântico.
    \item \textbf{Caminho de Expansão (Decodificador):} Realiza a amostragem ascendente (*Up-sampling* ou convolução transposta) dos mapas de características, restaurando gradativamente a resolução espacial original da imagem para viabilizar a predição pixel a pixel.
\end{enumerate}

O grande diferencial da U-Net reside nas \textit{connections de salto} (*skip connections*). Essas estruturas copiam os mapas de características de alta resolução do caminho de contração e os concatenam diretamente aos mapas correspondentes no caminho de expansão. Matematicamente, se $x_l$ representa a saída da camada $l$ do codificador e $y_l$ representa a entrada da camada de \textit{up-sampling} correspondente, a operação de fusão pode ser denotada por:

\begin{equation}
    z_l = [ \mathcal{W}_{up}(y_l) \parallel x_l ]
\end{equation}

Onde $\parallel$ representa o operador de concatenação ao longo do eixo dos canais e $\mathcal{W}_{up}$ denota a operação de convolução transposta. Esse mecanismo preserva informações espaciais de granulação fina (como bordas e texturas microestruturais), as quais costumam ser mitigadas pelas camadas de subamostragem.

\section{Estimativa de Movimento via Fluxo Óptico Denso de Farneback}
\label{sec:farneback}

O fluxo óptico descreve o padrão de movimento aparente de objetos, superfícies e bordas em uma sequência visual, gerado pelo deslocamento relativo entre o observador e a cena. Diferente das abordagens esparsas, o algoritmo de Farneback computa um fluxo óptico denso, estimando um vetor de deslocamento para todos os pixels de cada quadro \cite{farneback2003two}.

O método fundamenta-se na aproximação da vizinhança de cada pixel por uma expansão polinomial quadrática. O modelo de intensidade local de um \textit{frame} no instante temporal $t_1$ pode ser expresso matricialmente como uma forma quadrática:

\begin{equation}
    f_1(\mathbf{x}) \sim \mathbf{x}^T \mathbf{A}_1 \mathbf{x} + \mathbf{b}_1^T \mathbf{x} + c_1
\end{equation}

Onde $\mathbf{x} = [x, y]^T$ representa as coordenadas espaciais do pixel, $\mathbf{A}_1$ é uma matriz simétrica $2\times2$, $\mathbf{b}_1$ é um vetor coluna e $c_1$ é um escalar, obtidos por meio do ajuste de mínimos quadrados ponderados na vizinhança. Assumindo que um quadro subsequente no instante $t_2$ é gerado por um deslocamento global ideal $\mathbf{d}$, o novo polinômio assume a forma:

\begin{equation}
    f_2(\mathbf{x}) = f_1(\mathbf{x} - \mathbf{d}) \sim (\mathbf{x} - \mathbf{d})^T \mathbf{A}_1 (\mathbf{x} - \mathbf{d}) + \mathbf{b}_1^T (\mathbf{x} - \mathbf{d}) + c_1
\end{equation}

Expandindo os termos e igualando os coeficientes correspondentes aos de um novo polinômio estimado para o segundo quadro ($f_2(\mathbf{x}) \sim \mathbf{x}^T \mathbf{A}_2 \mathbf{x} + \mathbf{b}_2^T \mathbf{x} + c_2$), obtém-se o sistema de equações que define o deslocamento:

\begin{equation}
    \mathbf{A}_2 = \mathbf{A}_1
\end{equation}
\begin{equation}
    \mathbf{b}_2 = \mathbf{b}_1 - 2\mathbf{A}_1\mathbf{d}
\end{equation}

A partir de $\mathbf{b}_2 - \mathbf{b}_1 = -2\mathbf{A}_1\mathbf{d}$, o vetor de deslocamento $\mathbf{d}$ para cada pixel pode ser teoricamente isolado resolvendo-se o sistema linear:

\begin{equation}
    \mathbf{A}_1\mathbf{d} = -\frac{1}{2}(\mathbf{b}_2 - \mathbf{b}_1)
\end{equation}

Na prática, como a igualdade da matriz de curvatura $\mathbf{A}_2 = \mathbf{A}_1$ raramente é perfeita devido a ruídos de amostragem e variações de iluminação, adota-se uma aproximação média $\mathbf{A}(\mathbf{x}) = \frac{\mathbf{A}_1(\mathbf{x}) + \mathbf{A}_2(\mathbf{x})}{2}$ e define-se o vetor de variação $\Delta\mathbf{b}(\mathbf{x}) = -\frac{1}{2}(\mathbf{b}_2(\mathbf{x}) - \mathbf{b}_1(\mathbf{x}))$. A solução robusta é então computada integrando-se as restrições em uma vizinhança local $\mathcal{I}$:

\begin{equation}
    \mathbf{d}(\mathbf{x}) = \left( \sum_{\mathbf{y} \in \mathcal{I}} \mathbf{A}(\mathbf{y})^T \mathbf{A}(\mathbf{y}) \right)^{-1} \sum_{\mathbf{y} \in \mathcal{I}} \mathbf{A}(\mathbf{y})^T \Delta\mathbf{b}(\mathbf{y})
\end{equation}

\section{Modelagem de Segundo Plano por Mistura de Gaussianas (MOG2)}
\label{sec:mog2}

Para a geração automática das máscaras de \textit{Ground Truth} sob regime de supervisão fraca, emprega-se a segmentação adaptativa de primeiro plano baseada em modelos de Mistura de Gaussianas (*Gaussian Mixture Models* -- GMM). O algoritmo clássico de Zivkovic aprimora o método tradicional ao selecionar dinamicamente o número de componentes gaussianas necessárias para modelar cada pixel em tempo real, reduzindo o custo computacional e aumentando a fidelidade da segmentação \cite{zivkovic2004improved}.

A história visual de um pixel específico até o instante corrente $t$ é modelada como uma série temporal $X = \{\mathbf{x}_1, \dots, \mathbf{x}_t\}$. A densidade de probabilidade de observar o valor corrente do pixel $\mathbf{x}_t$ é definida como uma soma ponderada de $K$ distribuições normais multivariadas:

\begin{equation}
    P(\mathbf{x}_t) = \sum_{m=1}^{K} w_{m,t} \cdot \eta(\mathbf{x}_t; \boldsymbol{\mu}_{m,t}, \boldsymbol{\Sigma}_{m,t})
\end{equation}

Onde $w_{m,t}$ representa a estimativa do peso da $m$-ésima gaussiana no tempo $t$, $\boldsymbol{\mu}_{m,t}$ denota o vetor de médias e $\boldsymbol{\Sigma}_{m,t} = \sigma_{m,t}^2 \mathbf{I}$ denota a matriz de covariância, assumindo independência entre os canais de cor. A distribuição de probabilidade normal $\eta$ é descrita por:

\begin{equation}
    \eta(\mathbf{x}_t; \boldsymbol{\mu}, \boldsymbol{\Sigma}) = \frac{1}{(2\pi)^{n/2} |\boldsymbol{\Sigma}|^{1/2}} \exp\left( -\frac{1}{2}(\mathbf{x}_t - \boldsymbol{\mu})^T \boldsymbol{\Sigma}^{-1} (\mathbf{x}_t - \boldsymbol{\mu}) \right)
\end{equation}

As componentes são ordenadas decrescentemente com base na razão $\frac{w_m}{\sigma_m}$. As primeiras $B$ distribuições que satisfazem o limiar de acumulação $T$ são consideradas como o modelo estável do plano de fundo:

\begin{equation}
    B = \arg\min_b \left( \sum_{m=1}^{b} w_{m,t} > T \right)
\end{equation}

Se o valor do pixel corrente $\mathbf{x}_t$ não apresentar uma distância estatística aceitável (geralmente medida via distância de Mahalanobis dentro de um intervalo de $2.5\sigma$) em relação a nenhuma das componentes pertencentes a $B$, o pixel é classificado como \textit{foreground} (primeiro plano ou objeto em movimento).

\section{Equalização Adaptativa de Histograma Limitada por Contraste (CLAHE)}
\label{sec:clahe}

Em cenários de camuflagem crítica, superfícies de textura altamente homogênea degradam a capacidade de correspondência vetorial do algoritmo de Farneback, induzindo o erro conhecido como o \textit{problema da abertura}. O algoritmo CLAHE (*Contrast Limited Adaptive Histogram Equalization*) opera como um estágio essencial de pré-processamento para mitigar esse viés, maximizando o contraste local sem gerar saturações espúrias \cite{zuiderveld1994contrast}.

Diferente da equalização de histograma global, o CLAHE subdivide a imagem em pequenas sub-regiões contextuais não sobrepostas chamadas de \textit{tiles} (geralmente de dimensões $8\times8$ pixels). Para cada uma dessas células, calcula-se um histograma local independente. Para evitar a amplificação excessiva de ruídos de alta frequência em regiões homogêneas, o CLAHE introduz um teto de corte (\textit{clip limit}) no histograma.

Os pixels que excedem o limite estabelecido pelo \textit{clip limit} são cortados e uniformemente redistribuídos entre todas as classes (\textit{bins}) do histograma local antes da computação da Função de Distribuição Acumulada (CDF). Para eliminar descontinuidades visuais e artefatos de bloco nas fronteiras dos \textit{tiles}, os mapeamentos finais dos níveis de intensidade dos pixels são calculados através de uma interpolação bilinear contínua das funções de transformação das quatro células adjacentes.

\section{Métricas para Validação}
\label{sec:metricas_validacao}

Para avaliar rigorosamente o desempenho dos modelos de aprendizado profundo desenvolvidos neste trabalho, é imprescindível o uso de métricas estatísticas sólidas. Em cenários de desbalanceamento severo de classes, como o encontrado na detecção de pequenos insetos camuflados, métricas globais podem ser enganosas, justificando a adoção de indicadores específicos baseados na matriz de confusão. A matriz contabiliza os Verdadeiros Positivos (TP - \textit{True Positives}), Falsos Positivos (FP - \textit{False Positives}), Verdadeiros Negativos (TN - \textit{True Negatives}) e Falsos Negativos (FN - \textit{False Negatives}).

\subsection{Acurácia (Accuracy)}
A acurácia mensura a proporção geral de predições corretas em relação ao total de pixels avaliados. Matematicamente, é definida como:

\begin{equation}
    Accuracy = \frac{TP + TN}{TP + TN + FP + FN}
\end{equation}

Embora útil em contextos balanceados, a acurácia é insuficiente em cenários onde a classe majoritária (fundo foliar) domina mais de $99\%$ da imagem, pois classificar toda a imagem como fundo resultaria em uma falsa percepção de alta precisão.

\subsection{Precisão (Precision)}
A precisão quantifica a proporção de verdadeiros positivos dentre todas as predições classificadas como positivas pelo modelo. É fundamental para mensurar a confiabilidade da rede e minimizar os falsos alarmes causados por variações de textura da folha.

\begin{equation}
    Precision = \frac{TP}{TP + FP}
\end{equation}

\subsection{Sensibilidade (Recall ou Sensitivity)}
Também conhecida como Taxa de Verdadeiros Positivos (TPR), a sensibilidade mede a capacidade do modelo de encontrar todos os alvos reais. No contexto do Manejo Integrado de Pragas, é crucial garantir que nenhum inseto passe despercebido.

\begin{equation}
    Recall = \frac{TP}{TP + FN}
\end{equation}

\subsection{F1-Score}
O \textit{F1-Score} consiste na média harmônica entre a Precisão e a Sensibilidade. Fornece uma métrica escalar balanceada que penaliza valores discrepantes, sendo fortemente recomendada para avaliar a eficácia do modelo em detectar os espécimes camuflados.

\begin{equation}
    F1\text{-}Score = 2 \times \frac{Precision \times Recall}{Precision + Recall}
\end{equation}

\subsection{Interseção sobre a União (IoU)}
Também referida como Índice de Jaccard, a IoU é a métrica padrão-ouro na segmentação semântica. Ela calcula a razão entre a área de interseção e a área de união da máscara predita pelo modelo com a máscara de \textit{Ground Truth}.

\begin{equation}
    IoU = \frac{TP}{TP + FP + FN}
\end{equation}

Uma IoU de $1.0$ indica uma sobreposição perfeita entre a predição e o alvo real, enquanto valores próximos a $0$ indicam ausência de sobreposição ou uma predição difusa.

\section{Estratégias para Lidar com o Desbalanceamento de Classes}
\label{sec:desbalanceamento}

A detecção de alvos em pequena escala impõe um cenário probabilístico onde os gradientes de erro gerados pela classe minoritária são frequentemente anulados pela classe do plano de fundo durante o treinamento por retropropagação. Para mitigar esse efeito, duas estratégias complementares são abordadas na literatura:
\begin{enumerate}
    \item \textbf{Aumento Extensivo de Dados (\textit{Data Augmentation}):} Aplicação de transformações espaciais (rotação, corte, espelhamento) e de intensidade (ruído e variação de brilho) para forçar o modelo a aprender representações invariantes e aumentar artificialmente o volume de amostragem da classe minoritária.
    \item \textbf{Adaptação da Função de Perda:} A substituição da tradicional Entropia Cruzada por funções sensíveis à proporção de pixels, como a \textit{Dice Loss} ou a \textit{Focal Loss}, garante que erros em regiões diminutas penalizem a rede de forma muito mais agressiva do que erros cometidos na classe de fundo.
\end{enumerate}

\section{Trabalhos Relacionados}
\label{sec:trabalhos_relacionados}

O uso de algoritmos de visão computacional na agricultura tem apresentado resultados notáveis. Pesquisas recentes focadas em \textit{Deep Learning} relataram o êxito de CNNs clássicas (como VGG16 e ResNet) na identificação de doenças foliares a partir de lesões visualmente contrastantes. No entanto, a problemática da camuflagem e do mimetismo revela a fragilidade dos descritores puramente morfológicos.

No âmbito de segmentação de objetos camuflados (COS - \textit{Camouflaged Object Segmentation}), abordagens que adotam tensores multicanais para capturar sinais de movimento vêm se destacando. Trabalhos focados na fusão de fluxo óptico e dados RGB evidenciam que o rastreamento da assinatura cinemática não só quebra as barreiras da camuflagem cromática como também melhora a robustez do \textit{pipeline} analítico contra as oscilações dinâmicas de iluminação inerentes a um ambiente não controlado, pavimentando o caminho para o uso da U-Net estendida investigada nesta pesquisa.